% -*- coding: utf-8 -*-
% 问题三两阶段求解流程图（服务日入列 + 事件驱动时刻 + 覆盖兜底 + 逐日阶段二）
\begin{figure}[htbp]
\centering
\begin{tikzpicture}[
    box/.style={draw, rounded corners=3pt, minimum width=3.4cm, minimum height=0.75cm,
        align=center, font=\footnotesize, fill=blue!8},
    boxg/.style={draw, rounded corners=3pt, minimum width=3.4cm, minimum height=0.75cm,
        align=center, font=\footnotesize, fill=green!10},
    boxo/.style={draw, rounded corners=3pt, minimum width=3.4cm, minimum height=0.75cm,
        align=center, font=\footnotesize, fill=orange!12},
    boxr/.style={draw, rounded corners=3pt, minimum width=3.4cm, minimum height=0.75cm,
        align=center, font=\footnotesize, fill=red!8},
    boxp/.style={draw, rounded corners=3pt, minimum width=3.4cm, minimum height=0.75cm,
        align=center, font=\footnotesize, fill=purple!10},
    arrow/.style={->, >=stealth, thick, draw=gray!70},
    darrow/.style={->, >=stealth, thick, draw=gray!70, dashed},
]

% 阶段一标题
\node[font=\small\bfseries, fill=green!15, draw, rounded corners, minimum width=12.5cm, minimum height=0.5cm] (s1) at (0, 4.6) {阶段一：固定排班（3840人）};

% 第一行：预处理与列生成（4 列）
\node[boxg] (window) at (-4.6, 3.5) {时间窗解析\\可行日期集 $D_p$（服务日自由）};
\node[boxg] (cols) at (0, 3.5) {全周候选列池\\事件驱动时刻 $\times$ (机场,机型) $\times$ 多设施序列};
\node[boxg] (load) at (4.6, 3.5) {EDF 双序满员装载\\mandatory 优先 + 临时填残座};
\node[boxp] (cover) at (8.2, 3.5) {覆盖兜底\\专属直达列族};

% 第二行：剪枝与 MILP（3 列）
\node[box] (prune) at (-3.4, 2.2) {列池剪枝 \\(40,10) $\to$ 约 7 万列};
\node[box] (mip) at (1.2, 2.2) {阶段一 COPT MILP\\$\min T$，3840 全覆盖 + 机队并发};
\node[boxg] (assign1) at (5.8, 2.2) {飞机指派\\区间图并发保证};

% 阶段一输出
\node[boxg] (tstar) at (0, 1.0) {每日 $T^*_d$（日预算）$\Rightarrow T^*$};

% 阶段二标题
\node[font=\small\bfseries, fill=orange!15, draw, rounded corners, minimum width=12.5cm, minimum height=0.5cm] (s2) at (0, -0.6) {阶段二：临时任务加入（160人，$T \le T^*$）};

% 阶段二流程（3 列）
\node[boxo, below=0.35cm of s2, xshift=-3.2cm] (day_cols) {逐日列裁剪\\当日可行临时需求};
\node[boxo, below=0.35cm of s2] (prize) {逐日 prize-collecting MIP\\$\max \sum z_p$，日预算 $T^*_d$};
\node[boxo, below=0.35cm of s2, xshift=3.2cm] (rerun) {阶段一架次作保底列\\MIP start 防无解};

\node[boxr, below=0.55cm of prize] (output) {输出 q3-routes / q3-assignments\\（3840 全覆盖，临时按预算承运）};

% 主箭头
\draw[arrow] (window) -- (cols);
\draw[arrow] (cols) -- (load);
\draw[arrow] (load) -- (cover);
\draw[arrow] (cover.west) -- (prune.east);
\draw[arrow] (prune) -- (mip);
\draw[arrow] (mip) -- (assign1);
\draw[arrow] (assign1.south) |- (tstar.east);
\draw[arrow] (tstar) -- (s2);
\draw[arrow] (s2) -- (day_cols);
\draw[arrow] (s2) -- (prize);
\draw[arrow] (s2) -- (rerun);
\draw[arrow] (day_cols) -- (prize);
\draw[arrow] (rerun) -- (prize);
\draw[arrow] (prize) -- (output);

% 覆盖兜底反馈
\draw[darrow] (cover.south) -- ++(0,-0.15) -| (window.south);
\node[font=\scriptsize, align=center, anchor=north] at (8.2, 2.75) {缺失人员\\补列};

\end{tikzpicture}
\caption{问题三两阶段求解流程（服务日入列 + 事件驱动时刻 + 覆盖兜底 + 逐日阶段二）}
\label{fig:q3_flow}
\end{figure}
